\subsubsection{Drain Coordinator(s)}
\label{drain-coordinators}
The Drain Coordinator(s) shall:

\begin{enumerate}
 \item
  Recruit a suitable staff of volunteers.
 \item
  Train all employees at the beginning of their term.
 \item
  Develop creative ideas for new apparel and merchandise.
 \item
  Organize times for the sale of leather jackets, and help the companies in charge of these items with the sales process, including receiving any deliveries.
 \item
  Collect and record all payments from students for the purchase of leather jackets.
 \item
  Keep accurate and regular records of inventory.
 \item
  Make cash deposits into the MES safe on a daily basis so that the cash in the Drain does not exceed \$200 at the end of each day, to a maximum of \$50 in coins and \$150 in cash.
 \item
  Ensure that all deposits in the MES safe are properly labelled with the source of income, amount, and date deposited.
 \item
  Be responsible for delegating the opening and closing of the Drain.
 \item
  Arrange for the Advertising Committee to promote sales and other campaigns.
 \item
  Consult with suppliers, students, and the VPF about upcoming campaigns (i.e. new paraphernalia).
 \item
  Keep the Drain neat and organized.
 \item
  Coordinate the sale of event tickets by Drain staff with the DoE.
 \item
  Keep a record of all sales (both ticket and item) made, via regular inventory checks and consolidating the Square sales records. A record of ticket sales is to be kept on the event sheet provided by the DoE.
 \item
  Report to the VPF.
 \item
  Be a position held by a maximum of three people.
\end{enumerate}